% ==================== قسمت (الف) ====================
\subsectionaddtolist{الف}

برای سری فوریه‌ی مثلثاتی، رابطه‌ی میان نرخ کاهش ضرایب و هموار بودن تابع به صورت زیر است:
\begin{itemize}
  \item ضرایب از مرتبه‌ی $O(1/n)$: تابع مربع‌انتگرال‌پذیر، لزوماً پیوسته نیست.
  \item ضرایب از مرتبه‌ی $O(1/n^2)$: تابع پیوسته است ولی مشتق اول آن پیوسته نیست.
  \item ضرایب از مرتبه‌ی $O(1/n^3)$: تابع پیوسته و مشتق اول آن نیز پیوسته است ولی مشتق دوم پیوسته نیست.
\end{itemize}

برای ضرایب داده‌شده، رفتار مجانبی برای $n\to\infty$ داریم:
$$a_n = \frac{\pi n}{n^4+1} \sim \frac{\pi}{n^3} = O\!\left(\frac{1}{n^3}\right)$$
$$b_n = \frac{\pi}{n^2+\pi^2+n\pi} \sim \frac{\pi}{n^2} = O\!\left(\frac{1}{n^2}\right)$$

ضرایب $b_n$ که کندتر میل به صفر می‌کنند، رفتار کلی سری را تعیین می‌کنند؛ بنابراین ضرایب از مرتبه‌ی $O(1/n^2)$ هستند.

\medskip
پس تابع $f(t)$ {پیوسته} است، اما مشتق اول آن ممکن است در برخی نقاط گسستگی پرشی داشته باشد.

% ==================== قسمت (ب) ====================
\subsectionaddtolist{ب}

با تکرار ساده با دوره‌ی $T=1$، در مرزهای دوره داریم:
$$f(0^+) = 0,\quad f(1^-) = 1-1 = 0 \quad\Rightarrow\quad \text{تابع پیوسته است.}$$
اما مشتق:
$$f'(t)=1-2t \Rightarrow f'(0^+)=1,\quad f'(1^-)=-1$$
در مرز دوره مشتق پرش $\Delta f' = 1-(-1)=2$ دارد، بنابراین ضرایب از مرتبه‌ی $O(1/n^2)$ هستند.

برای پیوسته‌کردن مشتق در مرز، تابع را با بازتاب فرد نسبت به $t=1$ گسترش می‌دهیم:
$$g(t) = \begin{cases} t - t^2 & 0 < t < 1 \\ -(2-t) + (2-t)^2 = t^2 - 3t + 2 & 1 < t < 2 \end{cases}$$
و این تابع را با دوره‌ی $T=2$ تکرار می‌کنیم.

با گسترش فرد-تناوبی با دوره‌ی $T=2$ به صورت:
$$\boxed{g(t) = \begin{cases} t-t^2 & 0 < t < 1 \\ t^2-3t+2 & 1 < t < 2 \end{cases}, \quad g(t+2)=g(t)}$$
تابع $g(t)$ پیوسته و مشتقش نیز پیوسته است، در نتیجه ضرایب سری فوریه‌اش از مرتبه‌ی $O(1/n^3)$ به صفر همگرا می‌شوند.
